\section{Analisis de resultados}

\subsection{Medición de resistencia}

Al realizar un análisis comparativo entre los devanados Cuadro 7.3, se confirma que el de Alta Tensión presenta una resistencia significativamente mayor al de Baja Tensión. Este resultado es físicamente esperado, ya que, para elevar la tensión, el transformador requiere un mayor número de espiras empleando un conductor de menor calibre (sección transversal reducida), lo cual incrementa la longitud total del devanado y por consecuencia su resistencia eléctrica. Finalmente, la propagación de la incertidumbre instrumental durante este proceso ratifica que el método del voltímetro-amperímetro fue la técnica adecuada para el nivel de corriente manejado en el ensayo.

\subsection{Medición de relación de transformación}


\subsection{Verificación de polaridad}

La prueba de polaridad permitió determinar la orientación relativa de los devanados del transformador mediante la configuración del ensayo, donde se aplicó una tensión al devanado de alta y se unieron terminales adyacentes de ambos devanados, midiendo la tensión resultante entre los terminales restantes ($E_x = 185 \text{ V}$). Dado que $E_x > E_p$, el resultado confirma experimentalmente que el equipo presenta polaridad aditiva, cumpliendo con la condición:

\begin{equation}
    E_x \approx E_p + E_s
\end{equation}

Este concepto de polaridad aditiva indica una configuración interna específica donde los puntos de igual polaridad instantánea (homólogos) de los devanados de alta y baja tensión se encuentran situados de manera diagonalmente opuesta en la disposición física del transformador.


\subsection{Ensayo de relación de las pérdidas debido a las cargas y tensiones de cortocircuito}
\subsection{Ensayo de la medición de las pérdidas y la corriente de vacío}

